\documentclass[11pt]{article}
\usepackage[margin=1in]{geometry}
\usepackage{graphicx}
\usepackage{hyperref}
\hypersetup{colorlinks=true, linkcolor=blue, citecolor=blue, urlcolor=blue}
\usepackage{listings}
\usepackage{xcolor}
\usepackage{booktabs}
\usepackage{amsmath}
\usepackage[numbers]{natbib}

\lstset{
  basicstyle=\small\ttfamily,
  breaklines=true,
  frame=single,
  numbers=left,
  numberstyle=\tiny,
  columns=fullflexible
}


\title{KathAI Project Member Application}
\author{Your Name}
\date{\today}

\begin{document}
\maketitle

\section{General Questionnaire}

\subsection{Essentials}

\paragraph{1. Tell us something about yourself.}
%

\paragraph{2. What does RAG really do? What is semantic meaning according to your computer?}
%

\paragraph{3. Why is dealing with noisy data important?}
%

\paragraph{4. Why do we need RAG? Can't LLMs solve all our NLP problems?}
%

\subsection{Project Goals and Expectations}

\paragraph{1. Describe the goals of the project in as much detail as possible.}
%

\paragraph{2. Imagine yourself a year later, looking back at the work you did.}
%

\paragraph{3. Considering your definition of tremendous success, how much effort would it take?}
%

\subsection{Motivation and Drive}

\paragraph{1. What is your primary motivation for joining this project?}
%

\paragraph{2. Is there another opportunity more aligned with this motivation than KathAI PM?}
%

\paragraph{3. What could dampen your enthusiasm or make you disengage?}
%

\paragraph{4. Depth vs.\ breadth tradeoff, early UG — is this a good opportunity for you?}
%

\subsection{Commitments/PoRs}

\paragraph{1. Other commitments this academic year, peaks, and how you'll prioritize.}
%

\section{Technical Questionnaire}

\subsection{The Damaged Recording (2.1)}

\subsubsection*{Part 1: Speech Degradation \& Baseline}

\paragraph{1. Representations of the reference recording.}
%
% \includegraphics[width=\linewidth]{figures/2.1_representations.png}

\paragraph{2. Degraded versions (one factor at a time).}
%

\paragraph{3. Recoverable vs.\ unrecoverable degradation.}
%

\subsubsection*{Part 2: Reconstruction \& Failure Modes}

\paragraph{4. Two restoration techniques.}
%

\paragraph{5. Reference vs.\ damaged vs.\ restored comparison.}
%

\paragraph{6. Break the pipeline.}
%

\paragraph{7. ``Cleaner-sounding audio is not always higher-fidelity audio.''}
%

\subsection{The Chunk That Cut the Story in Half (2.2)}

\paragraph{1. Why naive fixed-size chunking breaks on raw transcripts.}
%

\paragraph{2. Two alternative chunking strategies.}
%

\paragraph{3. Implementation and evaluation of context loss.}
%
% \begin{table}[h]
% \centering
% \begin{tabular}{lcc}
% \toprule
% Strategy & Hit rate @ k & Notes \\
% \midrule
% Semantic & & \\
% Speaker-turn & & \\
% \bottomrule
% \end{tabular}
% \caption{Retrieval evaluation across chunking strategies}
% \end{table}

\paragraph{4. Bonus: Whisper transcripts made queryable.}
%

\subsection{When the English Sounds Right but the Meaning Is Wrong (2.3)}

\subsubsection*{Part 1: Where does the meaning go?}

\paragraph{1. Data-flow diagram through the translation pipeline.}
%
% \includegraphics[width=\linewidth]{figures/2.3_dataflow.png}

\paragraph{2. Behavior under meaningfully different linguistic situations.}
%

\subsubsection*{Part 2: Can you improve a failure?}

\paragraph{3. One repeated failure, diagnosis, and fix.}
%

\paragraph{4. Did the intervention actually improve translation?}
%

\paragraph{5. Flagging unreliable translations without ground truth.}
%

\section*{Resources}
% Cite everything from notes/citations.md here.
\bibliographystyle{plainnat}
\bibliography{references}

\end{document}
